% Exact safe statement/formula inputs; target proof prose remains embargoed.
% Input cdp:st:C.1.1; series Paper II; source SHA256 74497bfe327fae25c9cdcb17334b3bf42a9793a9df7c0a4404d97568f2f760f0.
\begin{proposition}\hypertarget{tag-01BN}{}\label{cdp:st:C.1.1}


Let $Y$ be a connected compact complex threefold with
$b_1(Y)=b_2(Y)=0$. Every holomorphic line bundle $L$ on $Y$ satisfies
\[
 \chi(Y,T_Y\otimes L)=\frac12\int_Yc_3(T_Y)
 =1-\frac{b_3(Y)}2.
\]
The orientation in the integral is the complex orientation. The group
$\Pic(Y)$ has finitely many connected components, each a torsor for the
complex vector group $H^1(Y,\OO_Y)$, and
$\pi_0\Pic(Y)\simeq H^2(Y;\Z)$ through the first Chern class.
\end{proposition}

% Input cdp:st:E7; series Paper II; source SHA256 74497bfe327fae25c9cdcb17334b3bf42a9793a9df7c0a4404d97568f2f760f0.
\begin{proposition}[Application to the named six-sphere families]\hypertarget{tag-01GW}{}\label{cdp:st:E7}


Every member of the following torus-fibred families,
including the period family $X_c$, its members with $|p|>1$,
Viaclovsky's $X_{3k,4l}$, his two-parameter family and his combs,
and the Liu--Rollenske family, satisfies the rank hypothesis of \ref{cdp:st:E6}.
Consequently every sufficiently small deformation retains a
holomorphic torus fibration over $\mathbb P^1$. Every such member
with $b_2=0$, and every sufficiently small deformation of it, has
algebraic dimension one. In particular the fibration-retention
question in Section~10(1) of the supplied comb paper has an
affirmative answer.
\end{proposition}

